%=======================================================================
%  CHAPTER 7 — NUMERICAL PROBLEMS AND SOLUTIONS
%=======================================================================
\section*{\color{acc}Ch 7 \textperiodcentered{} Multiple Access \gap
\textcolor{sub}{Numerical Problems \& Solutions}}
\vspace{-4pt}{\color{acc}\rule{\textwidth}{1.1pt}}\vspace{4pt}

{\color{sub}\footnotesize Every problem below is a real PYQ, with its year tags and marks,
except the four marked \pillo{sub}{PROBLEM}, which are tutorial or textbook problems added to
cover a pattern the papers imply but never set. Solutions are worked in full: what was
calculated, why, and which formula was used. \mk{Nothing is skipped to save space.}}

\lead{The four formulas everything here is built on}
\begin{itemize}
  \item \textbf{FDMA channels} \; $N = \dfrac{B_t - 2 B_{\text{guard}}}{B_c}$
  \item \textbf{TDMA channels} \; $N = \dfrac{m\,(B_{\text{tot}} - 2 B_{\text{guard}})}{B_c}$,
        \; $m =$ slots per carrier
  \item \textbf{TDMA frame efficiency} \;
        $\eta_f = \left(1 - \dfrac{b_{OH}}{b_T}\right) \times 100\,\%$
  \item \textbf{Hadamard recursion} \;
        $H_{2N} = \begin{bmatrix} H_N & H_N \\ H_N & -H_N \end{bmatrix}$,
        \; $H_1 = [\,1\,]$
\end{itemize}
{\color{sub}\footnotesize $B_t$ $=$ total spectrum allocation, $B_{\text{guard}}$ $=$ guard band
at \emph{one} edge, $B_c$ $=$ channel (carrier) bandwidth, $b_{OH}$ $=$ overhead bits per frame,
$b_T$ $=$ total bits per frame. \mk{All three bandwidths must be in the same unit before you
divide.} Convert everything to Hz, or everything to kHz, and never mix.}

\hr

\T{Problem 1 \quad Number of channels available in an FDMA system}
{\color{sub}\footnotesize \tF{4} \yr{\texttt{80 Ba} \m{3+2}, \textbf{79 Ch} \m{4},
\textbf{77 Ch} \m{4}, \textbf{78 Ch} \m{4}} \gap$\cdot$\gap
\mk{same formula three times}, only the numbers change}

\lead{1.1 \; $B_t = 25$ MHz, guard $= 100$ kHz, channel $= 200$ kHz \yr{(\texttt{80 Ba})}}
\begin{asked}{2080 Baishakh \textperiodcentered{} Q7(a) \textperiodcentered{} [3+2]}
Explain the non-linear effect in FDMA. If the total spectrum allocation is 25 MHz, the guard
band allocated at the edge of the spectrum is 100 KHz, and the channel bandwidth is 200 KHz.
Find the number of channels available in an FDMA system.
\par\vspace{1pt}
{\color{sub}The 3-mark theory half is answered in the chapter, \S7.2. Only the 2-mark
calculation is worked here.}
\end{asked}

\begin{itemize}
  \item \textbf{Step 1: list the given data and put it all in one unit.} The formula divides
        bandwidths by each other, so a mixed unit silently gives an answer 1000 times wrong.
  \begin{itemize}
    \item $B_t = 25 \text{ MHz} = 25 \times 10^6$ Hz
    \item $B_{\text{guard}} = 100 \text{ kHz} = 100 \times 10^3 = 0.1 \times 10^6$ Hz
    \item $B_c = 200 \text{ kHz} = 200 \times 10^3 = 0.2 \times 10^6$ Hz
  \end{itemize}
  \item \textbf{Step 2: subtract the guard bands.} \mk{A guard band sits at \emph{each} edge of
        the allocation}, at the bottom and at the top, so $2 B_{\text{guard}}$ is lost before
        any channel is cut. This factor of 2 is the single most common slip in this problem.
        \[ B_t - 2B_{\text{guard}} \;=\; 25 \times 10^6 - 2\,(100 \times 10^3)
           \;=\; 25{,}000{,}000 - 200{,}000 \;=\; 24{,}800{,}000 \text{ Hz} \]
  \item \textbf{Step 3: divide the usable spectrum by the channel bandwidth.}
        \[ N \;=\; \frac{B_t - 2B_{\text{guard}}}{B_c}
           \;=\; \frac{24{,}800{,}000}{200 \times 10^3}
           \;=\; \frac{24{,}800}{200} \;=\; 124 \]
  \item \textbf{Step 4: state the answer as a whole number.} $N$ counts channels, so it is an
        integer. Here the division is exact, and where it is not, \mk{round \emph{down}}: a
        fractional channel cannot be assigned.
  \item \mk{Final answer: $N = \mathbf{124}$ channels.}
\end{itemize}

\lead{1.2 \; $B_t = 12.8$ MHz, guard $= 10$ kHz, channel $= 30$ kHz
\yr{(\textbf{79 Ch}, \textbf{77 Ch})}}
\begin{asked}{2079 Chaitra \textperiodcentered{} Q8(b) \; and \; 2077 Chaitra \textperiodcentered{} Q10(b) \textperiodcentered{} [4]}
Calculate the number of radio channels available in the FDMA system with following data. A US
analog mobile phone system is allocated 12.8 MHz fir every simplex band. The total spectrum
allocated is 12.8 MHz, the guard bandwidth is 10 KHz and the channel bandwidth is 30 KHz.
\par\vspace{1pt}
{\color{sub}``fir'' is the 2079 Chaitra paper's own typo for ``for''; 2077 Chaitra prints the
same question correctly. The two papers are otherwise word for word identical.}
\end{asked}

\begin{itemize}
  \item \textbf{Step 1: read the statement carefully.} It says 12.8 MHz twice: once as ``allocated
        for every simplex band'' and once as ``the total spectrum allocated''. \mk{These are the
        same number, not two numbers to add.} $B_t = 12.8$ MHz is what goes into the formula.
  \begin{itemize}
    \item The phrase ``every simplex band'' is telling you this is \textbf{one direction only}.
          A full duplex system would need a second 12.8 MHz for the reverse link, and the
          channel count would still be 12.8 MHz worth.
  \end{itemize}
  \item \textbf{Step 2: convert to one unit.} Working in kHz keeps the numbers small:
        $B_t = 12{,}800$ kHz, \; $B_{\text{guard}} = 10$ kHz, \; $B_c = 30$ kHz.
  \item \textbf{Step 3: remove both guard bands.}
        \[ B_t - 2B_{\text{guard}} \;=\; 12{,}800 - 2(10) \;=\; 12{,}800 - 20
           \;=\; 12{,}780 \text{ kHz} \]
  \item \textbf{Step 4: divide by the channel bandwidth.}
        \[ N \;=\; \frac{12{,}780}{30} \;=\; 426 \]
  \item \mk{Final answer: $N = \mathbf{426}$ radio channels.}
\end{itemize}

\lead{1.3 \; $B_t = 12.5$ MHz, guard $= 10$ kHz, channel $= 30$ kHz \yr{(\textbf{78 Ch})}}
\begin{asked}{2078 Chaitra \textperiodcentered{} Q7(b) \textperiodcentered{} [4]}
If total bandwidth is 12.5 MHz, guard bandwidth is 10 KHz and channel bandwidth is 30 KHz,
find the number of channels available in an FDMA system.
\end{asked}

\begin{itemize}
  \item \textbf{Step 1: same formula, in kHz.}
        $B_t = 12{,}500$ kHz, \; $B_{\text{guard}} = 10$ kHz, \; $B_c = 30$ kHz.
  \item \textbf{Step 2: remove both guard bands.}
        \[ B_t - 2B_{\text{guard}} \;=\; 12{,}500 - 20 \;=\; 12{,}480 \text{ kHz} \]
  \item \textbf{Step 3: divide.}
        \[ N \;=\; \frac{12{,}480}{30} \;=\; 416 \]
  \item \mk{Final answer: $N = \mathbf{416}$ channels.}
  \item \textbf{Worth knowing:} 416 is the real AMPS figure. In the US, each cellular carrier
        was allocated 416 channels per simplex band. \mk{If your answer to this question is not
        416, you have made an arithmetic slip.}
\end{itemize}

\lead{1.4 \; The same problem run backwards \; \pillo{sub}{practice}}
\begin{asked}[PROBLEM]{not a PYQ \textperiodcentered{} AS Sir, Tutorial VIII, Example 2 (alternate form)}
In US AMPS, 416 channels are allocated to various cellular operators. The channel between them
is 30 kHz with the guard band of 10 kHz. Calculate the spectrum allocation given to each
operator.
\end{asked}

\begin{itemize}
  \item \textbf{Why this is worth doing:} the papers have only ever asked for $N$. The same
        formula rearranged asks for $B_t$, and \mk{it is the obvious next variant} for a paper
        to set.
  \item \textbf{Step 1: rearrange the channel-count formula for $B_t$.} Starting from
        $N = (B_t - 2B_{\text{guard}})/B_c$, multiply both sides by $B_c$ and add
        $2B_{\text{guard}}$:
        \[ B_t \;=\; N B_c + 2 B_{\text{guard}} \]
  \item \textbf{Step 2: substitute.}
        \[ B_t \;=\; 416 \times (30 \times 10^3) \;+\; 2\,(10 \times 10^3) \]
  \item \textbf{Step 3: evaluate each term.}
  \begin{itemize}
    \item Channels: $416 \times 30{,}000 = 12{,}480{,}000$ Hz $= 12.48$ MHz.
    \item Guard bands: $2 \times 10{,}000 = 20{,}000$ Hz $= 0.02$ MHz.
  \end{itemize}
  \item \textbf{Step 4: add.}
        \[ B_t \;=\; 12{,}480{,}000 + 20{,}000 \;=\; 12{,}500{,}000 \text{ Hz} \]
  \item \mk{Final answer: $B_t = \mathbf{12.5}$ MHz for each simplex band.}
  \item \textbf{Check:} this is exactly the input of Problem 1.3, which returned 416. The two
        problems are inverses of each other, so each one checks the other.
\end{itemize}

\hr

\T{Problem 2 \quad Simultaneously transmitting users in an N-TDMA cluster}
{\color{sub}\footnotesize \tP{1} \yr{\textbf{80 Ch} \m{4}}}

\begin{asked}{2080 Chaitra \textperiodcentered{} Q8(b) \textperiodcentered{} [4]}
In the N-TDMA system, it uses one way bandwidth of 25 MHz for the forward (or reverse) channel.
The system bandwidth is divided into radio channels of 30 kHz. Guard bands with
$B_g = 20$ kHz are used. Calculate number of simultaneously transmitting users that can be
accommodated in each cluster is.
\par\vspace{1pt}
{\color{sub}The trailing ``is'' is the paper's own slip. ``Cluster'' here means the whole
system bandwidth, since in a cluster every channel is used exactly once.}
\end{asked}

\begin{itemize}
  \item \textbf{Step 1: understand what is being counted.} A \textbf{cluster} is the group of
        cells over which the entire allocation is used once, with no repetition. So the number
        of simultaneously transmitting users per cluster is simply
        \mk{the number of channels the whole allocation supports}.
  \begin{itemize}
    \item Frequency reuse means the \emph{next} cluster reuses the same channels, so the answer
          is per cluster, not per system.
    \item \mk{No cluster size $N$ is needed}, and none is given. If the question had asked for
          users \emph{per cell}, you would then divide by $N$.
  \end{itemize}
  \item \textbf{Step 2: note which formula applies.} The statement gives no number of time slots
        per carrier, so take $m = 1$ and the TDMA formula collapses to the FDMA one:
        \[ N_u \;=\; \frac{m\,(B_s - 2B_g)}{B_c} \;\xrightarrow{\;m=1\;}\;
           \frac{B_s - 2B_g}{B_c} \]
  \item \textbf{Step 3: list the data in kHz.}
  \begin{itemize}
    \item $B_s = 25 \text{ MHz} = 25{,}000$ kHz (one-way, so no halving is needed: the
          statement already says forward \emph{or} reverse)
    \item $B_g = 20$ kHz \; \item $B_c = 30$ kHz
  \end{itemize}
  \item \textbf{Step 4: remove both guard bands and divide.}
        \[ N_u \;=\; \frac{25{,}000 - 2(20)}{30} \;=\; \frac{25{,}000 - 40}{30}
           \;=\; \frac{24{,}960}{30} \;=\; 832 \]
  \item \mk{Final answer: $N_u = \mathbf{832}$ simultaneously transmitting users per cluster.}
  \item \textbf{Sanity check:} the guard bands cost only $40/25{,}000 = 0.16\,\%$ of the
        allocation, so the answer must land just under $25{,}000/30 = 833.3$. It does.
\end{itemize}

\hr

\T{Problem 3 \quad TDMA frame efficiency in GSM}
{\color{sub}\footnotesize \tP{2} \yr{\textbf{70 Bh} \m{4}, 74 Ma \m{2}} \gap$\cdot$\gap
\mk{the numbers 156.25, 1250, 322 and 74.24 \% are worth memorising outright}}

\lead{3.1 \; Find the frame efficiency \yr{(\textbf{70 Bh})}}
\begin{asked}{2070 Bhadra \textperiodcentered{} Q8 \textperiodcentered{} [2+6+4]}
Define multiple access. What are the merits and demerits of Code Division Multiple Access? If a
normal GSM time slot consists of 6 trailing bits, 8.25 guard bits, 26 training bits, and
2 traffic bursts of 58 bits of data, find the frame efficiency.
\par\vspace{1pt}
{\color{sub}The first two parts are theory, answered in chapter \S7.1 and \S7.6. Only the
4-mark calculation is worked here.}
\end{asked}

\begin{itemize}
  \item \textbf{Step 1: build one time slot, field by field.} Add every field the statement
        lists. \mk{The 8.25 guard bits are not a typo}: GSM's guard period really is a quarter
        bit longer than 8 bits, so the slot length comes out non-integer.
\end{itemize}
\par\vspace{2pt}
\begin{tabularx}{\textwidth}{@{}L{4.2cm}L{3.0cm}L{2.4cm}X@{}}
\toprule
\hrow \thead{Field} & \thead{Count} & \thead{Bits} & \thead{Overhead or payload?} \\
\midrule
Trailing (tail) bits & $6$ & 6 & overhead, ramps the modulator \\
Guard bits & $8.25$ & 8.25 & overhead, absorbs propagation delay spread \\
Training bits & $26$ & 26 & overhead, trains the equalizer \\
Traffic bursts & $2 \times 58$ & 116 & \mk{payload}, the only useful bits \\
\midrule
\textbf{One time slot} & & \mk{$\mathbf{156.25}$} & \\
\bottomrule
\end{tabularx}

\begin{itemize}
  \item Arithmetic in full: \; $6 + 8.25 + 26 + 2(58) = 6 + 8.25 + 26 + 116 = 156.25$ bits.
  \item \textbf{Step 2: build the frame.} \mk{A GSM frame holds 8 time slots}, one per user.
        This number is not in the question, it is standard GSM knowledge you must supply.
        \[ b_T \;=\; 8 \times 156.25 \;=\; 1250 \text{ bits per frame} \]
  \item \textbf{Step 3: total the overhead over the whole frame.} Every one of the 8 slots
        carries its own tail, guard and training bits, so multiply each by 8:
        \[ b_{OH} \;=\; 8(6) \;+\; 8(8.25) \;+\; 8(26) \;=\; 48 + 66 + 208 \;=\; 322
           \text{ bits} \]
  \begin{itemize}
    \item $8 \times 6 = 48$ \; (tail) \; \item $8 \times 8.25 = 66$ \; (guard)
    \item $8 \times 26 = 208$ \; (training)
  \end{itemize}
  \item \textbf{Step 4: apply the frame efficiency formula.} Efficiency is the fraction of the
        frame that is \emph{not} overhead:
        \[ \eta_f \;=\; \left(1 - \frac{b_{OH}}{b_T}\right) \times 100\,\%
           \;=\; \left(1 - \frac{322}{1250}\right) \times 100\,\% \]
  \item \textbf{Step 5: evaluate.}
        \[ \frac{322}{1250} = 0.2576 \quad\rightarrow\quad
           1 - 0.2576 = 0.7424 \quad\rightarrow\quad \eta_f = 74.24\,\% \]
  \item \mk{Final answer: $\eta_f = \mathbf{74.24\,\%}$.}
  \item \textbf{Cross-check the other way.} Payload per frame $= 8 \times 116 = 928$ bits, and
        $928/1250 = 0.7424$. \mk{Same answer, which confirms that ``overhead'' was defined as
        everything except the two traffic bursts.}
\end{itemize}

\lead{3.2 \; Show that TDMA frame efficiency cannot reach 100 \% in GSM \yr{(74 Ma)}}
\begin{asked}{2074 Magh \textperiodcentered{} Q8 \textperiodcentered{} [6+2]}
Explain briefly channel structure of GSM. Show that TDMA frame efficiency cannot reach 100\% in
GSM
\par\vspace{1pt}
{\color{sub}The 6-mark GSM channel structure belongs to Ch 8. Only the 2-mark proof is here.}
\end{asked}

\begin{itemize}
  \item \textbf{Step 1: write the definition as an inequality, not a number.}
        \[ \eta_f = \left(1 - \frac{b_{OH}}{b_T}\right)\times 100\,\% \]
        so \; $\eta_f = 100\,\%$ \; requires \; $b_{OH} = 0$.
  \item \textbf{Step 2: show $b_{OH}$ can never be zero, field by field.} Each field is
        physically necessary, so none of them can be set to zero:
  \begin{itemize}
    \item \textbf{Guard bits} absorb the difference in propagation delay between a mobile at
          the cell edge and one beside the base station. Remove them and \mk{bursts from
          different mobiles overlap}, and both are lost.
    \item \textbf{Training bits} give the adaptive equalizer a known sequence to estimate the
          channel from. The channel changes between bursts, so this must be repeated
          \mk{in every burst}, not sent once.
    \item \textbf{Tail bits} let the modulator ramp up and down inside the slot. Without them
          the burst edges splatter energy into adjacent channels.
    \item \textbf{Preamble / synchronisation} bits let base station and subscriber identify
          each other and align in time.
  \end{itemize}
  \item \textbf{Step 3: conclude with the actual number.} In GSM,
        $b_{OH} = 322$ of $b_T = 1250$ bits, i.e. \mk{25.76 \% of every frame is overhead},
        giving $\eta_f = 74.24\,\%$. Since $b_{OH} = 322 > 0$, $\eta_f < 100\,\%$.
  \item \mk{Final answer: $\eta_f = 74.24\,\% < 100\,\%$, and it is bounded away from 100 \%
        because the burst structure that makes TDMA work is itself made of non-payload bits.}
  \item \textbf{The general statement}, worth one line: \mk{TDMA buys the ability to share a
        carrier by spending bits on framing.} FDMA transmits continuously and needs almost no
        framing, which is exactly why FDMA has lower overhead (chapter \S7.10).
\end{itemize}

\lead{3.3 \; GSM slot and frame timing, and simultaneous users \; \pillo{sub}{practice}}
\begin{asked}[PROBLEM]{not a PYQ \textperiodcentered{} AS Sir, Tutorial VIII, Examples 3 and 4}
Consider GSM, a TDMA/FDD system using 25 MHz for the forward link, broken into radio channels
of 200 kHz. If 8 speech channels are supported on a single radio channel and no guard band is
assumed, find the number of simultaneous users. If each frame consists of 8 time slots, each
slot contains 156.25 bits, and data is transmitted at 270.833 kbps, find (a) the time duration
of a bit, (b) of a slot, (c) of a frame, and (d) how long a user must wait between two
transmissions.
\end{asked}

\begin{itemize}
  \item \textbf{Users: step 1, find the bandwidth one user really occupies.} A 200 kHz carrier
        is shared by 8 users, so \mk{each user's share is $200/8 = 25$ kHz}, even though the
        user briefly uses the whole 200 kHz.
  \item \textbf{Users: step 2, divide the allocation by that share.} No guard band is assumed,
        so nothing is subtracted:
        \[ N \;=\; \frac{25 \text{ MHz}}{200 \text{ kHz} / 8}
           \;=\; \frac{25 \times 10^6}{25 \times 10^3} \;=\; 1000 \]
        \mk{GSM can accommodate 1000 simultaneous users} in 25 MHz.
  \item \textbf{(a) Time duration of a bit.} The bit period is the reciprocal of the bit rate:
        \[ T_b \;=\; \frac{1}{270.833 \text{ kbps}} \;=\; \frac{1}{270{,}833}
           \;=\; 3.692 \times 10^{-6} \text{ s} \;=\; \mathbf{3.692\ \mu s} \]
  \item \textbf{(b) Time duration of a slot.} A slot is 156.25 bits long, so multiply:
        \[ T_{\text{slot}} \;=\; 156.25 \times T_b \;=\; 156.25 \times 3.692\ \mu s
           \;=\; 576.9\ \mu s \;=\; \mathbf{0.577 \text{ ms}} \]
  \item \textbf{(c) Time duration of a frame.} A frame is 8 slots:
        \[ T_f \;=\; 8 \times T_{\text{slot}} \;=\; 8 \times 0.577 \text{ ms}
           \;=\; \mathbf{4.615 \text{ ms}} \]
  \begin{itemize}
    \item Cross-check through the bit count: $1250 \text{ bits} \times 3.692\ \mu s
          = 4.615$ ms. \checkmark
    \item So the frame rate is $1/4.615 \text{ ms} = 216.7$ frames per second, which also equals
          $270.833 \text{ kbps} / 1250 \text{ bits}$. \checkmark
  \end{itemize}
  \item \textbf{(d) Wait between transmissions.} A user owns one slot in every frame, so after
        its burst ends it must wait for the \emph{next} frame to come round:
        \mk{$\mathbf{4.615}$ ms}, one full frame time.
  \item \textbf{Why this matters:} 4.615 ms of buffering is why TDMA speech has more delay than
        FDMA speech, and the 0.577 ms burst is why the burst bit rate is 8 times the user's own
        rate, which is what forces the equalizer (chapter \S7.3).
\end{itemize}

\hr

\T{Problem 4 \quad Construct the $8 \times 8$ Hadamard code}
{\color{sub}\footnotesize \tF{3} \yr{\textbf{\texttt{82 Bh}} \m{3}, \textbf{\texttt{80 Bh}} \m{4},
\textbf{79 Ch} \m{3+5}}}

\begin{asked}{2079 Chaitra \textperiodcentered{} Q7 \m{3+5} \; \textperiodcentered{} \; also 2082 Bhadra Q6(a) \m{2+3}, 2080 Bhadra Q6(b) \m{4}}
Write the steps to satisfy condition for Hadamard code. Construct 8 by 8 Hadamard code that
satisfies above conditions.
\par\vspace{1pt}
{\color{sub}2082 Bhadra prints ``Define turbo coding. Construct 8 by 8 Hadamard code.''
and 2080 Bhadra prints ``Construct Hadamard matrix for $H_8$.'' Same construction, three
different mark splits.}
\end{asked}

\lead{Step 1 \; State the conditions the matrix must satisfy \m{3}}
\begin{itemize}
  \item It is a \textbf{square} matrix of order $N$, every entry $+1$ or $-1$.
  \item \mk{Every two distinct rows are orthogonal}: multiply them element by element, add up
        the products, and the result is 0.
  \item Equivalently $H_N H_N^{T} = N I_N$, where $I_N$ is the $N \times N$ identity.
  \item The order must be $N = 1$, $N = 2$, or a \textbf{multiple of 4}. There is no
        $3 \times 3$ or $6 \times 6$ Hadamard matrix.
  \item The $N$ rows are the $N$ \textbf{codewords}, each $N$ bits long, so the code carries
        $k = \log_2 N$ information bits. For $N = 8$: \mk{8 codewords, 8 bits each, $k = 3$}.
  \item Minimum distance is $d_{\min} = N/2 = 4$, so it detects 3 errors and corrects
        $\lfloor (4-1)/2 \rfloor = 1$ error.
\end{itemize}

\lead{Step 2 \; Start from $H_1$ and $H_2$}
\begin{itemize}
  \item The seed is the trivial $1 \times 1$ matrix \; $H_1 = \begin{bmatrix} 1 \end{bmatrix}$.
  \item Applying the recursion once gives the $2 \times 2$ matrix everybody starts from:
        \[ H_2 \;=\; \begin{bmatrix} H_1 & H_1 \\ H_1 & -H_1 \end{bmatrix}
           \;=\; \begin{bmatrix} 1 & 1 \\ 1 & -1 \end{bmatrix} \]
  \item \textbf{Check it:} row 1 times row 2 $= (1)(1) + (1)(-1) = 1 - 1 = 0$. \checkmark
        Orthogonal.
\end{itemize}

\lead{Step 3 \; Double it to $H_4$ using the Sylvester recursion}
\begin{itemize}
  \item The rule is \; $H_{2N} = \begin{bmatrix} H_N & H_N \\ H_N & -H_N \end{bmatrix}$:
        \mk{copy the block three times and negate the bottom right one}.
  \item $-H_2$ means flip the sign of every entry:
        $-H_2 = \begin{bmatrix} -1 & -1 \\ -1 & 1 \end{bmatrix}$.
  \item Assembling the four blocks:
        \[ H_4 \;=\; \begin{bmatrix} H_2 & H_2 \\ H_2 & -H_2 \end{bmatrix}
           \;=\;
           \begin{bmatrix}
             1 &  1 &  1 &  1 \\
             1 & -1 &  1 & -1 \\
             1 &  1 & -1 & -1 \\
             1 & -1 & -1 &  1
           \end{bmatrix} \]
  \item \textbf{Check a pair:} row 2 times row 3
        $= (1)(1) + (-1)(1) + (1)(-1) + (-1)(-1) = 1 - 1 - 1 + 1 = 0$. \checkmark
\end{itemize}

\lead{Step 4 \; Double once more to $H_8$ \m{5}}
\begin{itemize}
  \item Apply the same rule to $H_4$:
        \; $H_8 = \begin{bmatrix} H_4 & H_4 \\ H_4 & -H_4 \end{bmatrix}$.
  \item The top-left, top-right and bottom-left $4\times4$ blocks are all $H_4$ copied
        unchanged. The bottom-right block is $H_4$ with every sign flipped.
\end{itemize}
\[ H_8 \;=\;
\begin{bmatrix}
 1 &  1 &  1 &  1 &  1 &  1 &  1 &  1 \\
 1 & -1 &  1 & -1 &  1 & -1 &  1 & -1 \\
 1 &  1 & -1 & -1 &  1 &  1 & -1 & -1 \\
 1 & -1 & -1 &  1 &  1 & -1 & -1 &  1 \\
 1 &  1 &  1 &  1 & -1 & -1 & -1 & -1 \\
 1 & -1 &  1 & -1 & -1 &  1 & -1 &  1 \\
 1 &  1 & -1 & -1 & -1 & -1 &  1 &  1 \\
 1 & -1 & -1 &  1 & -1 &  1 &  1 & -1
\end{bmatrix} \]

\lead{Step 5 \; Verify the orthogonality condition}
\begin{itemize}
  \item Pick any two distinct rows, multiply element by element, add. \mk{Show at least one
        such check in the exam}, it is worth a mark.
  \item \textbf{Row 2 against row 3:}
        \[ (1)(1) + (-1)(1) + (1)(-1) + (-1)(-1) + (1)(1) + (-1)(1) + (1)(-1) + (-1)(-1) \]
        \[ = \; 1 - 1 - 1 + 1 + 1 - 1 - 1 + 1 \;=\; 0 \quad \checkmark \]
  \item \textbf{Row 1 against row 5:}
        \[ 1 + 1 + 1 + 1 - 1 - 1 - 1 - 1 \;=\; 0 \quad \checkmark \]
  \item \textbf{Any row against itself:}
        $1+1+1+1+1+1+1+1 = 8 = N$. \checkmark This is the second property CDMA decoding uses.
\end{itemize}

\lead{Step 6 \; Convert to binary Walsh codes}
\begin{itemize}
  \item For transmission the $\pm 1$ entries are mapped to bits: \mk{$+1 \rightarrow 0$ and
        $-1 \rightarrow 1$}.
  \item That turns the eight rows of $H_8$ into the eight \textbf{Walsh codes} $W_0 \dots W_7$:
\end{itemize}
\par\vspace{2pt}
\begin{tabularx}{\textwidth}{@{}L{1.9cm}L{3.3cm}XL{1.9cm}L{3.3cm}X@{}}
\toprule
\hrow \thead{Code} & \thead{Binary} & \thead{Ones} & \thead{Code} & \thead{Binary} & \thead{Ones} \\
\midrule
$W_0$ & \texttt{0 0 0 0 0 0 0 0} & 0 & $W_4$ & \texttt{0 0 0 0 1 1 1 1} & 4 \\
$W_1$ & \texttt{0 1 0 1 0 1 0 1} & 4 & $W_5$ & \texttt{0 1 0 1 1 0 1 0} & 4 \\
$W_2$ & \texttt{0 0 1 1 0 0 1 1} & 4 & $W_6$ & \texttt{0 0 1 1 1 1 0 0} & 4 \\
$W_3$ & \texttt{0 1 1 0 0 1 1 0} & 4 & $W_7$ & \texttt{0 1 1 0 1 0 0 1} & 4 \\
\bottomrule
\end{tabularx}
\begin{itemize}
  \item \textbf{Check:} every codeword except the all-zero one has exactly \mk{$N/2 = 4$ ones},
        which is the binary statement of $d_{\min} = N/2$.
  \item \textbf{Scaling up:} one more application of the recursion gives $H_{64}$, whose 64 rows
        are what IS-95 actually uses. $W_0$, the all-zero row, is its pilot channel. See
        chapter \S7.7 for where these codes sit in the CDMA transmitter.
\end{itemize}

\hr

\T{Problem 5 \quad CDMA encoding and decoding, worked with numbers}
{\color{sub}\footnotesize \tF{3} \yr{\textbf{\texttt{82 Bh}} \m{7}, \textbf{79 Ch} \m{8},
\textbf{77 Ch} \m{6}} \gap$\cdot$\gap
\mk{the paper always says ``with relevant example''}, so an answer with no arithmetic loses
most of the marks}

\begin{asked}{2082 Bhadra \textperiodcentered{} Q7(a) \m{7} \; \textperiodcentered{} \; 2079 Chaitra Q8(a) \m{8} \; \textperiodcentered{} \; 2077 Chaitra Q10(a) \m{2+6}}
Explain the implementation of CDMA with its encoding part and decoding part with relevant
example.
\par\vspace{1pt}
{\color{sub}2079 Chaitra prints ``Illustrate Encoding and decoding mechanism of CDMA system in
brief.'' 2077 Chaitra prints ``What are the major disadvantages of spread spectrum system?
Explain the implementation of CDMA system with appropriate example.'' The 2-mark
disadvantages rider is answered in chapter \S7.4 and \S7.6.}
\end{asked}

\lead{5.1 \; Four stations, all transmitting}
\begin{itemize}
  \item \textbf{Step 1: write down the rules the arithmetic uses.} State these before
        calculating, they are worth marks on their own.
  \begin{enumerate}
    \item Codes are \textbf{orthogonal}: $C_i \cdot C_j = 0$ for $i \ne j$.
    \item A code times itself gives $C_i \cdot C_i = N$, the number of stations. Here $N = 4$.
    \item \textbf{Data mapping}: bit $1 \rightarrow +1$, bit $0 \rightarrow -1$,
          silent $\rightarrow 0$.
    \item Multiplying a sequence by a number scales every element of it.
  \end{enumerate}
  \item \textbf{Step 2: set up the given data.} Four stations, data $d_1 = 1$, $d_2 = 0$,
        $d_3 = 1$, $d_4 = 1$, and the four orthogonal chip sequences:
\end{itemize}
\par\vspace{2pt}
\begin{tabularx}{\textwidth}{@{}L{2.4cm}L{4.4cm}L{2.4cm}X@{}}
\toprule
\hrow \thead{Station} & \thead{Chip sequence $C_i$} & \thead{Data bit $d_i$} &
\thead{Mapped value} \\
\midrule
ST1 & $C_1 = (+1,\,+1,\,+1,\,+1)$ & 1 & $+1$ \\
ST2 & $C_2 = (+1,\,-1,\,+1,\,-1)$ & 0 & $-1$ \\
ST3 & $C_3 = (+1,\,+1,\,-1,\,-1)$ & 1 & $+1$ \\
ST4 & $C_4 = (+1,\,-1,\,-1,\,+1)$ & 1 & $+1$ \\
\bottomrule
\end{tabularx}

\begin{itemize}
  \item \textbf{Step 3: confirm the codes really are orthogonal} before using them. Two spot
        checks are enough in the exam:
  \begin{itemize}
    \item $C_1 \cdot C_2 = (1)(1) + (1)(-1) + (1)(1) + (1)(-1) = 1 - 1 + 1 - 1 = 0$ \checkmark
    \item $C_3 \cdot C_4 = (1)(1) + (1)(-1) + (-1)(-1) + (-1)(1) = 1 - 1 + 1 - 1 = 0$ \checkmark
    \item $C_1 \cdot C_1 = 1 + 1 + 1 + 1 = 4 = N$ \checkmark
  \end{itemize}
  \item \textbf{Step 4: encode. Each station multiplies its own code by its own data value.}
  \begin{itemize}
    \item $d_1 C_1 = (+1)\,(+1,+1,+1,+1) \;=\; (+1,\,+1,\,+1,\,+1)$
    \item $d_2 C_2 = (-1)\,(+1,-1,+1,-1) \;=\; (-1,\,+1,\,-1,\,+1)$
    \item $d_3 C_3 = (+1)\,(+1,+1,-1,-1) \;=\; (+1,\,+1,\,-1,\,-1)$
    \item $d_4 C_4 = (+1)\,(+1,-1,-1,+1) \;=\; (+1,\,-1,\,-1,\,+1)$
  \end{itemize}
  \item \textbf{Step 5: all four transmit at once, and the channel adds them.} Sum the four
        sequences \textbf{element by element}, i.e. column by column:
\end{itemize}
\par\vspace{2pt}
\begin{tabularx}{\textwidth}{@{}L{3.4cm}XXXX@{}}
\toprule
\hrow \thead{Contribution} & \thead{chip 1} & \thead{chip 2} & \thead{chip 3} & \thead{chip 4} \\
\midrule
$d_1 C_1$ & $+1$ & $+1$ & $+1$ & $+1$ \\
$d_2 C_2$ & $-1$ & $+1$ & $-1$ & $+1$ \\
$d_3 C_3$ & $+1$ & $+1$ & $-1$ & $-1$ \\
$d_4 C_4$ & $+1$ & $-1$ & $-1$ & $+1$ \\
\midrule
\textbf{Common channel $S$} & \mk{$+2$} & \mk{$+2$} & \mk{$-2$} & \mk{$+2$} \\
\bottomrule
\end{tabularx}
\begin{itemize}
  \item Column by column: \; $1-1+1+1 = +2$, \; $1+1+1-1 = +2$, \; $1-1-1-1 = -2$, \;
        $1+1-1+1 = +2$.
  \item So the composite signal on the common channel is \; $S = (+2,\,+2,\,-2,\,+2)$.
  \item \textbf{Step 6: decode. Suppose ST3 wants ST2's data.} Multiply $S$ by ST2's code
        element by element, add all four products, then divide by $N = 4$:
        \[ S \times C_2 \;=\; (+2)(+1),\; (+2)(-1),\; (-2)(+1),\; (+2)(-1)
           \;=\; (+2,\,-2,\,-2,\,-2) \]
        \[ \text{sum} \;=\; 2 - 2 - 2 - 2 \;=\; -4
           \qquad\rightarrow\qquad \frac{-4}{4} \;=\; -1 \]
  \item Applying the mapping backwards, $-1$ means \mk{bit 0}, which is exactly $d_2$.
        \checkmark
  \item \textbf{Step 7: show the same receiver recovers every station.} This is the proof the
        scheme works, and it is where the last marks are:
\end{itemize}
\par\vspace{2pt}
\begin{tabularx}{\textwidth}{@{}L{2.0cm}L{4.6cm}L{2.0cm}L{1.8cm}X@{}}
\toprule
\hrow \thead{Decode with} & \thead{$S \times C_i$, element by element} & \thead{Sum} &
\thead{$\div\,4$} & \thead{Recovered bit} \\
\midrule
$C_1$ & $(+2,\,+2,\,-2,\,+2)$ & $+4$ & $+1$ & \mk{1} $= d_1$ \checkmark \\
$C_2$ & $(+2,\,-2,\,-2,\,-2)$ & $-4$ & $-1$ & \mk{0} $= d_2$ \checkmark \\
$C_3$ & $(+2,\,+2,\,+2,\,-2)$ & $+4$ & $+1$ & \mk{1} $= d_3$ \checkmark \\
$C_4$ & $(+2,\,-2,\,+2,\,+2)$ & $+4$ & $+1$ & \mk{1} $= d_4$ \checkmark \\
\bottomrule
\end{tabularx}
\begin{itemize}
  \item \mk{Final answer: all four data bits $(1,0,1,1)$ are recovered from one shared
        4-chip sequence}, with every station using the whole bandwidth for the whole time.
  \item \textbf{Why it works, in one line:} in $S \times C_k$ every cross term $d_i (C_i \cdot
        C_k)$ with $i \ne k$ is zero by orthogonality, so only $d_k (C_k \cdot C_k) = d_k N$
        survives. Dividing by $N$ returns $d_k$.
\end{itemize}

\lead{5.2 \; Three stations, one of them silent}
\begin{asked}[PROBLEM]{not a PYQ \textperiodcentered{} SST Ch 7, Figs 12.34 and 12.35 \textperiodcentered{} the textbook version of the same question}
Stations A, B and C are assigned codes $C_A = (+1,-1,-1,+1)$, $C_B = (+1,+1,+1,+1)$ and
$C_C = (+1,-1,+1,-1)$. Station A wants to transmit a 0 bit, station B wants to remain silent,
and station C wants to transmit a 1. Find the multiplexed CDMA signal, and show that a
demultiplexer recovers all three states.
\end{asked}

\begin{itemize}
  \item \textbf{Why work this one too:} it is the version drawn in the lecture figures, and it
        adds the \mk{silent station} case, which the four-station example does not cover. A
        paper asking ``with relevant example'' accepts either.
  \item \textbf{Step 1: map the three states.} Bit $0 \rightarrow -1$,
        \mk{silent $\rightarrow 0$}, bit $1 \rightarrow +1$.
  \begin{itemize}
    \item A: bit 0 $\rightarrow -1$ \; \item B: silent $\rightarrow 0$ \;
    \item C: bit 1 $\rightarrow +1$
  \end{itemize}
  \item \textbf{Step 2: encode each station.}
  \begin{itemize}
    \item $-1 \cdot C_A = (-1)(+1,-1,-1,+1) \;=\; (-1,\,+1,\,+1,\,-1)$
    \item $\phantom{-}0 \cdot C_B = (0)(+1,+1,+1,+1) \;=\; (0,\,0,\,0,\,0)$
          \; \mk{a silent station contributes nothing at all}
    \item $+1 \cdot C_C = (+1)(+1,-1,+1,-1) \;=\; (+1,\,-1,\,+1,\,-1)$
  \end{itemize}
  \item \textbf{Step 3: add element by element to get the multiplexed signal.}
  \begin{itemize}
    \item chip 1: $-1 + 0 + 1 = 0$ \; \item chip 2: $+1 + 0 - 1 = 0$
    \item chip 3: $+1 + 0 + 1 = +2$ \; \item chip 4: $-1 + 0 - 1 = -2$
  \end{itemize}
        \[ S \;=\; (0,\ 0,\ +2,\ -2) \]
  \item \textbf{Step 4: demultiplex each branch.} Multiply, add all bits, divide by $N = 4$:
\end{itemize}
\par\vspace{2pt}
\begin{tabularx}{\textwidth}{@{}L{2.0cm}L{4.6cm}L{2.0cm}L{1.8cm}X@{}}
\toprule
\hrow \thead{Branch} & \thead{$S \times C_i$, element by element} & \thead{Add all bits} &
\thead{$\div\,4$} & \thead{Result} \\
\midrule
A, $C_A$ & $(0,\,0,\,-2,\,-2)$ & $-4$ & $-1$ & \mk{bit 0} \checkmark \\
B, $C_B$ & $(0,\,0,\,+2,\,-2)$ & $0$ & $0$ & \mk{silent} \checkmark \\
C, $C_C$ & $(0,\,0,\,+2,\,+2)$ & $+4$ & $+1$ & \mk{bit 1} \checkmark \\
\bottomrule
\end{tabularx}
\begin{itemize}
  \item \mk{Final answer: $S = (0,0,+2,-2)$, and the decoder returns bit 0, silent, bit 1},
        matching what the three stations sent.
  \item \textbf{Note the three-way result.} The decoder does not just distinguish 1 from 0. It
        also reports \emph{silence} as an exact zero, which is what lets CDMA gain capacity from
        \mk{voice activity}: a silent user costs the system nothing.
\end{itemize}

\hr

\T{Problem 6 \quad Capacity of an IS-95 CDMA cell \; \pillo{sub}{practice}}
{\color{sub}\footnotesize \mk{No paper 2070--2082 has set this}, but it is the one calculation
the CDMA theory implies and it is a standard tutorial problem. Included so the pattern is
covered.}

\begin{asked}[PROBLEM]{not a PYQ \textperiodcentered{} AS Sir, Tutorial VIII, Example 8}
In an IS-95 CDMA system, if $W = 1.25$ MHz, $R = 9600$ bps and $N = 14$ users,
(a) calculate $E_b/N_0$;
(b) when there is no voice activity, calculate it for omnidirectional antennas;
(c) if voice activity $= 3/8$ and 3-sector antennas are used, calculate the total number of
users per cell.
\end{asked}

\begin{itemize}
  \item \textbf{Step 1: find the processing gain.} It is the ratio of spread bandwidth to data
        rate:
        \[ \frac{W}{R} \;=\; \frac{1.25 \times 10^6}{9600} \;=\; 130.21 \]
  \item \textbf{(a) Step 2: apply the single-cell CDMA capacity relation.} Each of the other
        $N-1$ users contributes an equal share of interference, so
        \[ \frac{E_b}{N_0} \;=\; \frac{W/R}{N - 1} \;=\; \frac{130.21}{14 - 1}
           \;=\; \frac{130.21}{13} \;=\; 10.02 \]
  \item In decibels, $10\log_{10}(10.02) = 10.01$ dB. \mk{Final answer:
        $E_b/N_0 \approx 10$, i.e. about 10 dB.}
  \begin{itemize}
    \item {\color{sub}The linear value 10.02 and the dB value 10.01 coincide here only because
          $10\log_{10} 10 = 10$. Do not assume that in other problems.}
  \end{itemize}
  \item \textbf{(b) With no voice activity and omnidirectional antennas, nothing changes.}
        Voice activity factor $\alpha = 1$ means every user talks all the time, which is the
        assumption already built into part (a). An omnidirectional antenna gives no
        sectorisation gain. \mk{So $E_b/N_0$ stays at 10.02, about 10 dB.}
  \item \textbf{(c) Step 3: rearrange the relation for $N$ and add the voice activity factor.}
        Solving $E_b/N_0 = (W/R)/(N-1)$ for $N$ gives $N = 1 + (W/R)/(E_b/N_0)$. A user who is
        silent a fraction of the time contributes interference only while talking, so the term
        is divided by $\alpha$:
        \[ N \;=\; 1 \;+\; \frac{1}{\alpha}\left[\frac{W/R}{E_b/N_0}\right] \]
  \item \textbf{Step 4: substitute} $\alpha = 3/8$, so $1/\alpha = 8/3$:
        \[ N \;=\; 1 + \frac{8}{3}\left[\frac{130.21}{10}\right]
           \;=\; 1 + \frac{8}{3}(13.021) \;=\; 1 + 34.72 \;=\; 35.72 \]
        i.e. about \textbf{35.7 users per sector}.
  \item \textbf{Step 5: apply the sectorisation gain.} A 3-sector antenna splits the cell into
        three, and each sector only hears the interference of its own users, so the cell carries
        three times the per-sector figure:
        \[ N_{\text{cell}} \;=\; 3 \times 35.72 \;=\; 107.2 \]
  \item \mk{Final answer: about $\mathbf{107}$ users per cell.}
  \item \textbf{Reading the result:} the raw processing gain alone would allow 14 users. Voice
        activity multiplies that by $8/3$ and sectoring by 3, taking it to roughly 107.
        \mk{This is exactly the ``soft capacity'' argument of chapter \S7.6 in numbers}: CDMA
        capacity is set by the interference budget, not by a fixed count of channels.
\end{itemize}
